\section{Isolated wells and deterministic energy tuning}
\label{sec:wells}

This section supplies an exponentially accurate classical procedure for
tuning the individual well energies. The procedure operates on
one-electron operators only; it does not diagonalize any of the spin or
Hubbard Hamiltonians from \Cref{sec:finite}.

There are two separate tasks. First we must show that moving a secondary
positive nucleus changes one isolated energy monotonically, with much
weaker effects on the other sites. Then we must evaluate those energies
accurately enough to implement the resulting contraction. The required
accuracy is exponential in $D$: an inverse-polynomial energy error in
these dilated units would grow after the final amplification. The
analytic and arithmetic estimates below address that precision cost.

Fix centers \(X_1,\ldots,X_m\in\R^3\) satisfying
\begin{equation}
 |X_i-X_j|\ge D-2S\quad(i\ne j),\qquad S\le D/100.
 \label{eq:wells-separation}
\end{equation}
Here \(D\) is a sufficiently large integer and \(m/D\) is smaller than a
sufficiently small absolute constant. All constants in this section are
uniform in the centers and the parameters \(\rho_i\in[1,4]\).
The stronger scale requirements used later imply these hypotheses.
Computational statements additionally assume rational centers. Their
bit lengths, rather than their numerical magnitudes, enter the
complexity bounds.

Choose an integer \(Z_*\) with \(e^D/2\le Z_*\le2e^D\). Round the following
two numbers to nearest integers, once and for all, before tuning:
\begin{equation}
\begin{split}
 Z_*q&=\operatorname{round}\!\left(
              Z_*\left(1-\frac{500m}{D}\right)\right),\\
 Z_*q'&=\operatorname{round}\!\left(Z_*\frac{1000m}{D}\right).
\end{split}
 \label{eq:wells-charges}
\end{equation}
Both integers are positive for the scales considered here. In particular,
\begin{equation}
\begin{split}
 q&=1-\frac{500m}{D}+O(Z_*^{-1}),\\
 q'&=\frac{1000m}{D}+O(Z_*^{-1}),\qquad
 c\frac mD\le q'\le C\frac mD.
\end{split}
 \label{eq:wells-charge-bounds}
\end{equation}
The primary and secondary nuclei of site \(i\) are at \(X_i\) and
\(X_i+\rho_i e_z\), respectively, in scaled coordinates. Write
\begin{equation}
\begin{split}
 r_i&=|y-X_i|,\qquad
 V_i(y)=\frac q{r_i}+\frac{q'}{|y-X_i-\rho_i e_z|},\\
 h&=-\frac12\Delta-\sum_{i=1}^m V_i.
\end{split}
 \label{eq:wells-full-one-body}
\end{equation}
These are the one-electron operators obtained from physical coordinates
by \(y=Z_*x\) and division of energies by \(Z_*^2\).

Fix a nondecreasing \(C^2\) function \(\chi:[0,\infty)\to[0,1]\) that
vanishes on \([0,1/16]\) and equals one on \([1/8,\infty)\). It can have
rational piecewise-polynomial data: integrate a normalized positive
multiple of \((t-1/16)^2(1/8-t)^2\) on the transition interval and extend
constantly outside. Set
\begin{equation}
 \chi_i(y)=\chi(r_i/D),\qquad
 W_i=\sum_{k\ne i}\chi_k V_k,\qquad
 h_i=-\frac12\Delta-V_i-W_i.
 \label{eq:wells-capped}
\end{equation}
The caps remove the other nuclear singularities from \(h_i\) while
retaining their fields near site \(i\).

\subsection{Uniform ground-state estimates}

We use the three-dimensional Sobolev bounds
\(\|u\|_6\le C\|u\|_{H^1}\) and
\(\|u\|_\infty\le C\|u\|_{H^2}\), together with their localized versions.
We also use the translated Hardy inequality
\begin{equation}
 \left\|\frac{u}{|\,\cdot-a\,|}\right\|_2
 \le2\|\nabla u\|_2\qquad(a\in\R^3).
 \label{eq:wells-hardy}
\end{equation}
Indeed, integrate
\(\operatorname{div}((y-a)/|y-a|^2)=|y-a|^{-2}\)
against \(|u|^2\) for a smooth compactly supported \(u\). Integration by
parts and Cauchy--Schwarz give
\[
 \|u/|\,\cdot-a\,|\|_2^2
 \le2\|u/|\,\cdot-a\,|\|_2\|\nabla u\|_2.
\]
Approximation proves \eqref{eq:wells-hardy} on \(H^1\).
Thus Coulomb multiplication maps \(H^1\) continuously into \(L^2\),
uniformly in the location of the pole.

\begin{proposition}[Uniform isolated-well estimates]
\label{prop:isolated-well}
The operators \(h_i\) are self-adjoint on \(H^2(\R^3)\), with form domain
\(H^1(\R^3)\). Each has a simple ground eigenvalue \(\lambda_i\) and a
real, positive, normalized eigenfunction \(\psi_i\). There is a constant
\(g>0\) such that the remaining spectrum is at least \(\lambda_i+g\).
For \(\phi_i(y)=\pi^{-1/2}e^{-r_i}\),
\begin{equation}
 |\lambda_i+1/2|\le C\frac mD,\qquad
 \|\psi_i-\phi_i\|_{H^2}\le C\frac mD.
 \label{eq:wells-h2}
\end{equation}
Moreover, for derivative orders \(j=0,1,2\),
\begin{equation}
 \|D^jW_i\|_\infty\le C\frac m{D^{j+1}},
 \qquad 0\le W_i\le\frac{100m}{D}.
 \label{eq:wells-field-bounds}
\end{equation}
These assertions hold before tuning and uniformly throughout
\([1,4]^m\).
\end{proposition}

\begin{proof}
On the nonzero support of \(\chi_k\), \(r_k\ge D/16\). Since \(\rho_k\le4\),
\[
 |y-X_k-\rho_k e_z|\ge r_k-4\ge r_k/2
\]
for large \(D\). The secondary singularity is inside the region where
the cap is identically zero. Differentiation through order two,
including derivatives of the cap, gives
\(\|D^j(\chi_kV_k)\|_\infty\le CD^{-j-1}\).
Summing proves the derivative bounds. For example, \(q\le1\), \(q'\le1\),
and sufficiently large \(D\) give a bound below \(100/D\) for each capped
potential, proving the stated function bound.

Put \(H_i^0=-\Delta/2-r_i^{-1}\) and \(B_i=h_i-H_i^0\).
Hardy's inequality and \eqref{eq:wells-field-bounds} give
\begin{equation}
 \|B_i u\|_2\le C\frac mD\|u\|_{H^1}.
 \label{eq:wells-perturbation}
\end{equation}
For any potential in these single-well operators, Hardy and Fourier
interpolation give, for each \(a>0\), a uniform bound
\(\|Vu\|_2\le a\|\Delta u\|_2+C_a\|u\|_2\).
The Kato--Rellich theorem for relatively bounded perturbations
\cite[Theorem~6.4, p.~135]{Teschl2009} therefore gives domain \(H^2\). The same estimates, with the Laplacian term
absorbed, prove uniform graph-norm equivalence:
\begin{equation}
 C^{-1}\|u\|_{H^2}
 \le \|h_i u\|_2+C\|u\|_2
 \le C\|u\|_{H^2}.
 \label{eq:wells-graph-norm}
\end{equation}
The corresponding statement holds for \(H_i^0\).
At form level a fixed positive shift of either operator is bounded
above and below by constant multiples of \(\|u\|_{H^1}^2\), giving
the asserted form domains.

The hydrogenic ground state and its isolation follow without an
excited-state formula. Direct integration by parts gives
\begin{equation}
 \langle u,(H_i^0+1/2)u\rangle
 =\frac12\left\|\nabla u+\frac{y-X_i}{r_i}u\right\|_2^2.
 \label{eq:wells-hydrogen-square}
\end{equation}
Its zero vectors are exactly the multiples of \(\phi_i\).
If normalized vectors orthogonal to \(\phi_i\) had energies tending to
\(-1/2\), they would be bounded in \(H^1\). The Coulomb form is compact
under weak convergence of such a sequence: within radius \(a\) of the
pole its absolute value is bounded by \(Ca\|u\|_{H^1}^2\), by Sobolev and
H\"older; beyond radius \(R\) it is bounded by \(R^{-1}\|u\|_2^2\);
on the intervening annulus local compactness applies. A weak limit
would therefore obey
\[
 \langle u,H_i^0u\rangle\le-1/2,\qquad
 \langle u,H_i^0u\rangle\ge-\tfrac12\|u\|_2^2\ge-1/2.
\]
Equality forces norm one and a multiple of \(\phi_i\), contrary to
orthogonality. This proves a fixed gap \(g_0>0\).

By \eqref{eq:wells-perturbation}, after a fixed positive shift the
perturbed form lies between \(1-Cm/D\) and \(1+Cm/D\) times the
unperturbed shifted form. The variational principle puts its first
level within \(Cm/D\) of \(-1/2\), and its second variational level at
least \(-1/2+g_0-Cm/D\). The negative first level is attained.
The preceding truncation argument gives form compactness of both own
Coulomb potentials and of the bounded decaying \(W_i\). A normalized
minimizing sequence is bounded in \(H^1\). Passing to a weakly convergent
subsequence with limit \(u\), compactness and lower semicontinuity give
\[
 \lambda_i\|u\|_2^2\le\langle u,h_i u\rangle\le\lambda_i<0.
\]
Since \(\|u\|_2\le1\), these inequalities force \(\|u\|_2=1\)
and attainment. For small enough \(m/D\) the
first level is simple and separated from the rest by a fixed \(g>0\).
Taking the absolute value of a real minimizer cannot increase its
energy. The Euler equation and the strong maximum principle on the
connected complement of the poles make it positive there.

The form bounds give \(\|\psi_i\|_{H^1}\le C\). Its equation implies
\[
 (H_i^0+1/2)\psi_i
   =(\lambda_i+1/2)\psi_i-B_i\psi_i,
\]
with right side of norm at most \(Cm/D\). On \(\phi_i^\perp\), the inverse
of \(H_i^0+1/2\) is bounded from \(L^2\) to \(H^2\): use the gap for the
\(L^2\) bound, then \eqref{eq:wells-graph-norm}.
The orthogonal component of \(\psi_i\) consequently has \(H^2\) norm
\(O(m/D)\). Normalization and the positive overlap with \(\phi_i\) put
its coefficient along \(\phi_i\) within \(O((m/D)^2)\) of one.
This proves \eqref{eq:wells-h2} and the uniform sup bound. Positivity
also holds at the poles by continuity and this sup-norm approximation,
since both poles lie within distance four of \(X_i\).
\end{proof}

\subsection{Moving the secondary nucleus}

The derivative of a moving Coulomb potential is not an \(L^2\)-bounded
multiplier. We instead control its expectations by Hardy's inequality.

\begin{lemma}[Bracketing and uniform secant bounds]
\label{lem:wells-secants}
With centers and rounded charges fixed, for every choice of the other
parameters,
\(\lambda_i(\rho_i=1)<-1/2<\lambda_i(\rho_i=4)\).
For \(1\le a<b\le4\), holding the other parameters fixed,
\begin{equation}
 c q'(b-a)\le\lambda_i(b)-\lambda_i(a)\le Cq'(b-a).
 \label{eq:wells-own-slope}
\end{equation}
If instead \(\rho_i\) is held fixed, then
\begin{equation}
 |\lambda_i(\rho)-\lambda_i(\widetilde\rho)|
 \le Cq'\frac m{D^2}\|\rho-\widetilde\rho\|_\infty.
 \label{eq:wells-other-slope}
\end{equation}
\end{lemma}

\begin{proof}
Taking the inner product of the eigen-equation with \(\phi_i\) gives
\begin{equation}
 \lambda_i=-\frac12+(1-q)-q'f(\rho_i)
       -\langle\phi_i,W_i\phi_i\rangle+O((m/D)^2),
 \label{eq:wells-energy-expansion}
\end{equation}
where
\[
 f(\rho)=\int_{\R^3}\frac{\pi^{-1}e^{-2|y|}}{|y-\rho e_z|}\dd y.
\]
Here \(\langle\phi_i,\psi_i\rangle=1+O((m/D)^2)\), and replacing
\(\psi_i\) by \(\phi_i\) in a perturbing expectation costs \(O((m/D)^2)\),
by \eqref{eq:wells-perturbation} and \eqref{eq:wells-h2}.
Also \(\langle\phi_i,r_i^{-1}\phi_i\rangle=1\).
The spherical average of \(|r\omega-\rho e_z|^{-1}\) is
\(1/\max(r,\rho)\), as direct integration in \(\omega\cdot e_z\) shows.
Consequently,
\begin{align}
 f(\rho)&=\rho^{-1}\int_0^\rho4r^2e^{-2r}\dd r
             +\int_\rho^\infty4r e^{-2r}\dd r,
 \label{eq:wells-radial-potential}\\
 f'(\rho)&=-\rho^{-2}\int_0^\rho4r^2e^{-2r}\dd r.
 \label{eq:wells-f-derivative}
\end{align}
In particular,
\[
 f(1)=1-2e^{-2}>0.7,\quad f(4)\le\tfrac14,\quad
 \frac{1-5e^{-2}}{16}\le-f'(\rho)\le1\quad(1\le\rho\le4).
\]
At the first endpoint the leading shift in
\eqref{eq:wells-energy-expansion} is at most \((500-700)m/D\);
at the second it is at least \((500-250-100)m/D\).
The quadratic and rounding errors fit inside these strict margins,
proving bracketing.

For the secant bound, fix a ground function \(u\) at either endpoint,
and put \(A(t)=X_i+t e_z\). The kernel derivative is dominated by
\(|y-A(t)|^{-2}\), and Hardy gives
\begin{align}
 \int\frac{|u^2-\phi_i^2|}{|y-A(t)|^2}\dd y
 &\le\left\|\frac{u-\phi_i}{|\,\cdot-A(t)\,|}\right\|_2
       \left\|\frac{u+\phi_i}{|\,\cdot-A(t)\,|}\right\|_2 \notag\\
 &\le4\|\nabla(u-\phi_i)\|_2
       \bigl(\|\nabla u\|_2+\|\nabla\phi_i\|_2\bigr)
 \le C\frac mD.
 \label{eq:wells-density-hardy}
\end{align}
The same bound with \(u^2\) justifies absolute integration in \(y\)
and \(t\). Thus replacing the density by \(\phi_i^2\) when integrating
the moving-pole kernel derivative incurs at most
\(Cq'(m/D)(b-a)\) error.
The variational principle gives
\[
 \langle\psi_i(b),(h_i(b)-h_i(a))\psi_i(b)\rangle
 \le\lambda_i(b)-\lambda_i(a)
 \le\langle\psi_i(a),(h_i(b)-h_i(a))\psi_i(a)\rangle.
\]
Both bounding quantities differ from \(-q'(f(b)-f(a))\) by the error
just estimated, proving \eqref{eq:wells-own-slope}.

For \(k\ne i\), the moving secondary pole in \(\chi_kV_k\) stays at
least \(D/32\) from the region where its cap is nonzero.
Its bounded potential changes in sup norm by at most
\(Cq'D^{-2}|\rho_k-\widetilde\rho_k|\).
Summing and applying the variational principle for a bounded additive
perturbation proves \eqref{eq:wells-other-slope}.
\end{proof}

\Needspace{10\baselineskip}
\subsection{Uniform tails}

\begin{lemma}[Crude tails before tuning]
\label{lem:well-tail}
Uniformly for all \(\rho\in[1,4]^m\),
\begin{equation}
 |\psi_i(y)|\le C\exp[-.99(r_i-.01D)],
 \label{eq:wells-tail}
\end{equation}
and, for sufficiently large \(D\),
\begin{equation}
 \|\psi_i\|_{H^1(\{r_i>.65D\})}\le e^{-.61D}.
 \label{eq:wells-h1-tail}
\end{equation}
\end{lemma}

\begin{proof}
On \(r_i\ge .01D\), both own poles are far away and
\(V_i+W_i=O(m/D)\). For \(a=.99\) and
\(b(y)=C\exp[-a(r_i-.01D)]\), direct calculation gives
\[
 \frac{(h_i-\lambda_i)b}{b}
 =-\frac{a^2}{2}+\frac a{r_i}-V_i-W_i-\lambda_i.
\]
Its constant part is \((1-a^2)/2+O(m/D)>0\) at the chosen scales.
The zero-order coefficient of \(h_i-\lambda_i\) is also positive
throughout this exterior region. The uniform sup bound fixes \(C\)
so that \(b\ge|\psi_i|\) at the inner boundary. An \(H^2(\R^3)\) function
is uniformly continuous and vanishes at infinity; otherwise uniform
continuity on a sequence of disjoint balls would contradict its
square integrability. On a large finite annulus add a constant
\(\varepsilon>0\) to \(b\) to dominate both signs of \(\psi_i\) at the
outer boundary. The supersolution inequality is preserved.
Testing the negative part of \(b+\varepsilon\mp\psi_i\) proves
comparison. Send the outer radius to infinity and then
\(\varepsilon\) to zero. Inside \(.01D\), the same sup bound proves
\eqref{eq:wells-tail}.

For the gradient estimate choose a smooth cutoff \(\zeta\) that vanishes
on \(r_i\le .645D\), equals one on \(r_i\ge .65D\), and has gradient
\(O(D^{-1})\). Testing the eigen-equation with \(\zeta^2\psi_i\) gives
\[
 \frac12\|\nabla(\zeta\psi_i)\|_2^2
 =\int(V_i+W_i+\lambda_i)|\zeta\psi_i|^2\dd y
       +\frac12\int|\nabla\zeta|^2|\psi_i|^2\dd y.
\]
All coefficients on this support are bounded. Thus
\eqref{eq:wells-tail} bounds the exterior \(H^1\) norm by a polynomial in
\(D\) times \(\exp[-.99(.645-.01)D]\).
Since \(.99(.645-.01)=.62865>.61\), the polynomial factor is absorbed
for large \(D\).
\end{proof}

\subsection{An analytic estimate with uniform constants}

We give the local regularity argument used in the numerical procedure.
It explains why the shrinking length scales near a moving Coulomb pole
do not increase the required polynomial degree beyond \(O(D)\).

\begin{lemma}[Uniform analytic patches]
\label{lem:wells-analytic}
Suppose \(u\) is bounded on the unit ball in \(\R^3\) and satisfies
\(\Delta u=Fu\) there, where
\begin{equation}
 \|\partial^\beta F\|_\infty\le A^{|\beta|+1}|\beta|!
 \quad\text{for every multi-index }\beta.
 \label{eq:wells-analytic-coefficient}
\end{equation}
There is a constant \(A_1\), depending only on \(A\), such that
\begin{equation}
 \max_{|\alpha|=k}\|\partial^\alpha u\|_{L^\infty(B_{1/4})}
 \le \|u\|_{L^2(B_1)}A_1^{k+1}k!.
 \label{eq:wells-analytic-derivatives}
\end{equation}
On a sufficiently small fixed ball about the origin, Taylor
polynomials consequently have geometrically decreasing errors in
both value and first derivative, with constants depending only on \(A\).
\end{lemma}

\begin{proof}
For two concentric balls of radii between \(1/2\) and \(1\), separated
by \(h\), the local estimates needed below are
\begin{align}
 \|D^2v\|_{2,\mathrm{inner}}
 &\le C\bigl(\|\Delta v\|_{2,\mathrm{outer}}
                   +h^{-2}\|v\|_{2,\mathrm{outer}}\bigr),
 \label{eq:wells-local-h2}\\
 \|\nabla v\|_{2,\mathrm{inner}}
 &\le C\bigl(h\|\Delta v\|_{2,\mathrm{outer}}
                   +h^{-1}\|v\|_{2,\mathrm{outer}}\bigr).
 \label{eq:wells-local-h1}
\end{align}
To obtain them, insert smooth cutoffs with first and second derivatives
\(O(h^{-1})\) and \(O(h^{-2})\). Integration by parts against a squared
cutoff and Cauchy--Schwarz, with \(ab\le (h a)^2/2+(b/h)^2/2\),
bound first derivatives on an intermediate ball as in
\eqref{eq:wells-local-h1}. Apply the whole-space Fourier inequality
\(\|D^2w\|_2\le C\|\Delta w\|_2\) to the cutoff function and substitute
that first-derivative bound in its commutator terms to obtain
\eqref{eq:wells-local-h2}. Approximation proves these estimates for
weak solutions. Successive localized Laplacian regularity makes \(u\)
smooth when \(F\) is smooth, justifying the following differentiations.

Fix an integer \(\ell\ge2\), set \(r_j=1-j/(2\ell)\), and write
\[
 T_j=\max_{|\alpha|=j}\|\partial^\alpha u\|_{L^2(B_{r_j})},
 \qquad M=\|u\|_{L^2(B_1)}.
\]
The first-derivative bound gives \(T_1\le C(\ell+A/\ell)M\).
For \(j\le\ell-2\), apply \eqref{eq:wells-local-h2} to an order-\(j\)
derivative between \(B_{r_j}\) and \(B_{r_{j+2}}\).
Leibniz' rule yields
\begin{equation}
 T_{j+2}\le C\left(
  \sum_{b=0}^j\binom jb A^{b+1}b!\,T_{j-b}
  +\ell^2T_j\right).
 \label{eq:wells-analytic-recurrence}
\end{equation}
Lower-order derivatives are available on the larger balls
\(B_{r_{j-b}}\). The multi-index coefficient sum of order \(b\)
is \(\binom jb\), by the multinomial identity.
If preceding orders satisfy \(T_s\le M(K\ell)^s\), the sum is at most
\[
 AM(K\ell)^j\sum_{b=0}^j(A/K)^b
 \le2AM(K\ell)^j
\]
for \(K\ge2A\), using \(\binom jb b!\le\ell^b\).
Increasing the fixed \(K\) absorbs the remaining \(C\ell^2\) factor
into \((K\ell)^2\). Both parity chains therefore satisfy
\(T_j\le M(K\ell)^j\).

Set \(\ell=k+2\). Sobolev embedding on fixed smaller balls bounds the
sup norm of an order-\(k\) derivative on \(B_{1/4}\) by
\(CM[K(k+2)]^{k+2}\).
The inequality \(k!\ge(k/e)^k\), with the small cases absorbed in
the constant, converts this into \eqref{eq:wells-analytic-derivatives}.
Taylor's formula, summed over multi-indices, bounds the degree-\(L\)
remainder on \(|z|\le a\) by a constant times \((C_1a)^{L+1}\).
The first-derivative remainder has an additional factor polynomial
in \(L\) and the same geometric decay. Fixing \(a\) small enough that
\(C_1a<1\) proves the last assertion.
\end{proof}

\subsection{A polynomial-size graded Ritz space}

Graded meshes with increasing polynomial degree are a standard way to
combine singular local behavior with exponential approximation away from
the singularities. For related eigenvalue approximation results with
singular potentials, see Maday and Marcati~\cite{MadayMarcati2019}.
Their bounded-domain discontinuous-Galerkin estimates are not an oracle
for the present moving attractive Coulomb wells. We construct the needed
conforming space and then certify its matrix arithmetic separately.
More specifically, the products of tensor-linear hats with local
polynomials below form a partition-of-unity approximation space, as in
Melenk and Babu\v{s}ka~\cite[Section~2]{MelenkBabuska1996}.
We give the cutoff and derivative estimates explicitly because the tensor
cells can have arbitrarily large aspect ratios.

This is a form-domain approximation: the piecewise polynomial
functions need only belong to \(H^1\). Their kinetic matrix uses first
weak derivatives, not second derivatives across element boundaries.

\begin{lemma}[Exponentially accurate Ritz approximation]
\label{lem:wells-ritz-approximation}
For each \(i\), each rational choice of centers and secondary parameters,
and each sufficiently large \(D\), one can explicitly construct a space
\(\mathcal V_i\subset H^1(\R^3)\) of rational piecewise polynomials,
of dimension polynomial in \(D\), supported in \(r_i<.8D\), such that
\begin{equation}
 \inf_{v\in\mathcal V_i}\|v-\psi_i\|_{H^1}\le e^{-.58D}.
 \label{eq:wells-ritz-h1}
\end{equation}
Its lowest Ritz value \(\lambda_i^{\mathrm R}\) satisfies
\begin{equation}
 0\le\lambda_i^{\mathrm R}-\lambda_i\le C e^{-1.16D}.
 \label{eq:wells-ritz-bias}
\end{equation}
The construction and constants are uniform over rational parameters
throughout \([1,4]^m\).
\end{lemma}

\begin{proof}
Work in local coordinates \(z=y-X_i\) on \([-D,D]^3\), with poles
at \(0\) and \(\rho_i e_z\). Put \(\delta=2^{-8D}\) and fix a sufficiently
small rational \(\xi>0\), as specified below. On each coordinate axis
let \(B\) be the set of projections of the two poles. Starting from
\([-D,D]\), bisect dyadically until every interval \(I\) obeys
\begin{equation}
 |I|\le\xi\min\{1,\dist(I,B)+\delta\}.
 \label{eq:wells-grid-rule}
\end{equation}
Every leaf length is at least \(\xi\delta/2\), since its parent failed
the stopping rule. Above a fixed spacing there are \(O(D)\) intervals.
At a finer level of length \(h\), only intervals within \(h/\xi\) of
one of the at most two projections can still split, so a fixed number
of intervals per level are split there.
There are \(O(\log(D/\delta))=O(D+\log D)\) levels, hence \(O(D)\)
intervals per axis. This remains true for poles arbitrarily close to
dyadic boundaries. The tensor product has \(O(D^3)\) nodes and cells.

Let \(b_w\) be the tensor-linear hat at node \(w\). The full hats form
a partition of unity in the cube. Retain nodes satisfying
\begin{equation}
 |w|<.73D,\qquad |w|>\delta,\qquad
 |w-\rho_i e_z|>\delta,
 \label{eq:wells-retained-nodes}
\end{equation}
and take \(b_w(z)(z-w)^\alpha\) for all \(|\alpha|\le L_0\), where \(L_0\)
is a sufficiently large fixed multiple of \(D\). These functions span
\(\mathcal V_i\), whose dimension before removing dependencies is
\(O(D^6)\).

Here is the support estimate that avoids any aspect-ratio assumption.
For either pole \(a\), every coordinate interval adjacent to \(w\) obeys
\[
 |I|\le\xi\min\{1,|w_j-a_j|+\delta\}.
\]
The support radius \(R_w^{\mathrm{hat}}\) of a retained hat thus satisfies
\begin{equation}
\begin{split}
 R_w^{\mathrm{hat}}
 &\le\xi\bigl(|w-a|+\sqrt3\delta\bigr)
   \le(1+\sqrt3)\xi|w-a|,\\
 R_w^{\mathrm{hat}}&\le\sqrt3\xi.
\end{split}
 \label{eq:wells-patch-radius}
\end{equation}
For \(d_w=\min\{1,|w|,|w-\rho_i e_z|\}\) this gives
\(R_w^{\mathrm{hat}}\le C\xi d_w\). Each support avoids both poles and,
for large \(D\), lies in \(|z|<.8D\).

Choose a fixed \(c<1/10\) and set \(R=c d_w\).
The whole ball \(B(w,R)\) avoids both poles. Throughout it every remote
primary has distance at least \(D-2S-.73D-c>D/8\), so it also avoids
every remote cap transition. On this enlarged ball the eigen-equation
contains only pure Coulomb potentials. Rescaling by \(R\) gives
\(\Delta u=Fu\) on the unit ball. Its coefficient satisfies
\eqref{eq:wells-analytic-coefficient} with an absolute \(A\):
local inverse-distance expansions give factorial derivative bounds
because \(R/\dist(w,a)\le c\); the extra \(R^2\) in the coefficient
makes \(R^2/\dist(w,a)\) bounded; and remote contributions sum to at
most \(Cm/D\). The eigenvalue term is bounded too.
The sup norm of \(u\) is uniform by \Cref{prop:isolated-well}.

Choose \(\xi\) sufficiently small compared to \(c\) and the fixed
analytic radius in \Cref{lem:wells-analytic}, putting the hat supports
inside that smaller ball after rescaling. With \(L_0\) a sufficiently
large fixed multiple of \(D\), its Taylor polynomial \(p_w\) satisfies
\begin{equation}
 |\psi_i-p_w|+R|\nabla\psi_i-\nabla p_w|\le e^{-20D}
 \quad\text{on }\supp b_w.
 \label{eq:wells-taylor-error}
\end{equation}
Increasing the fixed multiple absorbs the polynomial factor in the
differentiated remainder. The Taylor coefficients establish existence
of an approximant; the algorithm will not compute them.

To remove tiny neighborhoods of the poles and the exterior, let
\(a_{\mathrm{out}}\) equal one up to radius \(.66D\), decrease linearly
to zero by \(.72D\), and vanish thereafter. Let \(a_{\mathrm{in}}\) vanish
up to \(\delta\), increase linearly to one by \(2\delta\), and equal one
thereafter. Sample
\[
 a_{\mathrm{out}}(|z|)\,
 a_{\mathrm{in}}(|z|)\,
 a_{\mathrm{in}}(|z-\rho_i e_z|)
\]
at all grid nodes to get weights \(c_w\in[0,1]\), and put
\(b=\sum_w c_wb_w\). Nonzero weights occur only at retained nodes.
The function \(b\) equals one away from \(O(\delta)\)-balls about the
poles and away from \(|z|>.65D\): a node within \(2\delta\) of a pole
has adjacent lengths \(O(\xi\delta)\), whereas all intervals have
length at most \(\xi\). These facts localize the inner and outer
transitions, respectively.

Each first derivative of a tensor-linear interpolant is a convex
combination of edge secants in that coordinate. Consequently
\(|\nabla b|\le C/\delta\) on the inner transitions and
\(|\nabla b|\le C/D\) on the outer transition, irrespective of aspect
ratios. The bounded sup norm of \(\psi_i\) makes the inner
cutoff-gradient cost \(C\delta^{1/2}\). Also, by its \(H^2\) bound,
\[
 \|\nabla\psi_i\|_{L^2(B_{C\delta}(a))}
 \le C\delta\|\nabla\psi_i\|_6\le C\delta.
\]
The exterior cost follows from \eqref{eq:wells-h1-tail}. Thus
\begin{equation}
 \|\psi_i-b\psi_i\|_{H^1}
 \le C\bigl(\sqrt\delta+e^{-.61D}\bigr).
 \label{eq:wells-cutoff-error}
\end{equation}
The cutoff weights, like the Taylor coefficients, are only existence
coefficients in the explicitly generated space.

Set \(v=\sum_w c_wb_wp_w\in\mathcal V_i\).
At almost every point at most eight hats are nonzero. Their derivatives
cost at most \(C/\delta\), by the minimum interval length, and
\(R^{-1}\le C/\delta\).
On a region of volume \(O(D^3)\), \eqref{eq:wells-taylor-error} gives
\[
 \|b\psi_i-v\|_{H^1}\le C D^{3/2}\delta^{-1}e^{-20D}.
\]
Since \(\sqrt\delta=e^{-4(\log2)D}\) and
\(\delta^{-1}e^{-20D}=e^{-(20-8\log2)D}\),
this and \eqref{eq:wells-cutoff-error} prove
\eqref{eq:wells-ritz-h1} for large \(D\).

Finally, the form of \(h_i-\lambda_i\) is bounded in absolute value
by a constant times the \(H^1\) norm squared. For \(e=v-\psi_i\),
the eigen-equation cancels the cross terms, giving
\[
 \frac{\langle v,h_iv\rangle}{\|v\|_2^2}-\lambda_i
 =\frac{\langle e,(h_i-\lambda_i)e\rangle}{\|v\|_2^2}
 \le C\|e\|_{H^1}^2.
\]
Here \(\|v\|_2\) is bounded away from zero.
The variational principle supplies the opposite inequality for the
lowest Ritz value, proving \eqref{eq:wells-ritz-bias}.
\end{proof}

\subsection{Certified polynomial-bit eigenvalue evaluation}

A small Ritz space by itself does not establish bit complexity.
We next bound the conditioning loss and give certified matrix
entries and eigenvalue enclosures. All arithmetic below is classical.

\begin{proposition}[Uniform single-well eigenvalue oracle]
\label{prop:ritz-oracle}
Suppose the centers and the parameters in \([1,4]^m\) are rational,
and all their numerators and denominators, as well as those of
\(q,q'\), have at most \(B\) bits. Under the hypotheses of this section,
a deterministic algorithm returns rational numbers
\(\ell_i\le\lambda_i\le u_i\) with
\begin{equation}
 u_i-\ell_i\le\varepsilon_D,\qquad
 \varepsilon_D=2^{-\lceil8D/5\rceil}<e^{-1.1D}.
 \label{eq:wells-oracle-precision}
\end{equation}
Its bit complexity is polynomial in \(D,m,B\).
The bound is uniform over every rational sample in \([1,4]^m\);
no spectral gap of a Ritz matrix is assumed.
\end{proposition}

\begin{proof}
We divide the construction into exact polynomial arithmetic,
certified Coulomb integration, and rational eigenvalue enclosure.

\smallskip\noindent
\emph{Exact matrices and a mass lower bound.}
Construct the basis in \Cref{lem:wells-ritz-approximation}.
All grid decisions are rational comparisons: one-dimensional distances
in \eqref{eq:wells-grid-rule} are rational, and node membership in
\eqref{eq:wells-retained-nodes} is checked by comparing squared
distances. Grid depth is \(O(D+\log D)\). Endpoints, hat coefficients,
and coefficients of the degree-\(O(D)\) monomials have polynomial
bit length. For instance, raising a rational coordinate of \(b\)
bits to degree \(L\) costs at most \(O(Lb)\) bits, before polynomial
sums; here \(b=O(B+D+\log D)\) and \(L=O(D)\).
Only polynomially many coefficients and rational operations occur
in any polynomial expansion or box integral. Summing rational terms
by a product of their denominators gives a conservative polynomial
bit bound even without exploiting common denominators.

Let \(f_1,\ldots,f_{r_0}\) be this possibly redundant basis,
\(r_0=O(D^6)\). The mass and kinetic matrices
\begin{equation}
 M_{ab}=\int f_af_b\dd y,\qquad
 T_{ab}=\frac12\int\nabla f_a\cdot\nabla f_b\dd y
 \label{eq:wells-exact-matrices}
\end{equation}
are exact rationals obtained by box integration. A product of two
hats has only a fixed number of box pieces; alternatively, summing
over the entire polynomial-size grid also suffices.
Remove redundant functions by exact rational Gram-matrix elimination,
retaining an independent subset. This gives a positive definite
matrix \(M\) of dimension \(r\le r_0\), with corresponding kinetic
matrix \(T\). The subset is nonempty by the approximation property.

Let \(B_0\) be the largest bit length of a numerator or positive
denominator of the entries of this reduced \(M\), enlarged to at
least one. It is computed from the exact entries and is polynomial
in \(D,m,B\). If \(Q\) is the product of all entry denominators,
then \(Q\le2^{r^2B_0}\) and \(QM\) is an integer positive definite
matrix. Therefore \(\det(QM)\ge1\), while
\(\lambda_{\max}(M)\le r2^{B_0}\). The explicitly computable number
\begin{equation}
 \mu=2^{-K_M},\qquad
 K_M=r^3B_0+(r-1)\bigl(B_0+\lceil\log_2r\rceil\bigr)
 \label{eq:wells-mass-lower}
\end{equation}
satisfies \(\lambda_{\min}(M)\ge\mu\).
This intentionally coarse determinant bound is enough: \(K_M\)
is polynomial. Exact rank selection is itself polynomial-bit
arithmetic. After clearing denominators, fraction-free elimination
has entries given by minors of an integer matrix. Hadamard's
determinant bound bounds their bit lengths by a polynomial in its
dimension and entry bit length. Equivalently, ordinary rational
elimination with fractions reduced at each step has the same
polynomial bound through the bordered-minor formulas.

\smallskip\noindent
\emph{Certified Coulomb entries.}
The preceding determinant bound shows how many extra bits compensate
for a poorly conditioned basis. It remains to compute each potential
entry to that absolute accuracy. The graded supports avoid every pole,
which permits a uniformly convergent kernel expansion on each box.
On the supports of the basis functions every other cap equals one,
so every potential entry is a finite sum of pure Coulomb integrals
against polynomials on rational boxes. For each such box \(Q_0\),
each pole \(a\), and its midpoint \(c_0\), put
\[
 s=|c_0-a|^2,\qquad
 z(y)=\frac{2(c_0-a)\cdot(y-c_0)+|y-c_0|^2}{s}.
\]
Every contributing box is in the support of a retained hat.
By \eqref{eq:wells-patch-radius}, its diameter is a fixed small
fraction of its distance from either own pole. Remote poles are
farther still. Decreasing the fixed \(\xi\) if necessary gives
\(|z(y)|\le1/2\) throughout every such box. Thus
\begin{equation}
 \frac1{|y-a|}
 =s^{-1/2}\sum_{j=0}^J\binom{-1/2}{j}z(y)^j
       +\mathcal R_J(y),\qquad
 |\mathcal R_J(y)|\le s^{-1/2}2^{-J}.
 \label{eq:wells-coulomb-series}
\end{equation}
The coefficient bound
\(\left|\binom{-1/2}{j}\right|
 =4^{-j}\binom{2j}{j}\le1\)
proves this geometric remainder bound by summing the tail.

Here \(s\) and all coefficients of \(z\) are exact rationals.
Both \(\log_2^+(s^{-1/2})\) and the bit lengths of these rationals are
polynomial in \(D,m,B\); indeed the own-pole distances on supports
are bounded below by a fixed multiple of \(\delta\).
The positive scalar \(s^{-1/2}\) can be approximated to any requested
polynomial number of bits by rational bisection and exact squaring.

We spell out an absolute-error certificate, so cancellation among
polynomial terms causes no difficulty. If \(P\) is the polynomial
multiplying a kernel on \(Q_0\), a rational coefficient-and-volume
bound gives
\[
 A_{Q_0}\ge\int_{Q_0}|P(y)|\dd y.
\]
One may take the volume times the sum of absolute coefficients,
each multiplied by the appropriate coordinate bound to its degree.
Its bit length is polynomial. Include the charge in this bound.
If \(s^{-1/2}\le2^b\), the truncation error is at most
\(A_{Q_0}2^{b-J}\). An approximation to \(s^{-1/2}\) of absolute
error \(2^{-q_0}\) adds at most \(2A_{Q_0}2^{-q_0}\), since the
absolute sum in \eqref{eq:wells-coulomb-series} is at most two.
For \(T_0\) box--pole summands, taking
\begin{align*}
 J&\ge p+b+\lceil\log_2^+ A_{Q_0}\rceil
                         +\lceil\log_2 T_0\rceil+3,\\
 q_0&\ge p+\lceil\log_2^+ A_{Q_0}\rceil
                         +\lceil\log_2 T_0\rceil+4
\end{align*}
makes their total error smaller than \(2^{-p}\).
Use maxima over the summands to choose common orders.
Here \(\log_2^+t=\max\{0,\log_2t\}\); integer upper bounds can be
read directly from rational numerator and denominator bit lengths.

The polynomial \(z(y)^j\) has degree \(2j\) in three variables.
Even after multiplication by the basis polynomials it has only
\(O((D+J+1)^3)\) possible monomials. Polynomial multiplication,
integration, and coefficient bit lengths are polynomial in
\(D,m,B,J,q_0\). These bounds establish a deterministic entry
algorithm with any prescribed polynomial number \(p\) of absolute
precision bits. No matrix element is computed by integrating
across a Coulomb singularity or a cap transition.

\smallskip\noindent
\emph{Matrix error and exact eigenvalue decisions.}
Let \(A\) be the exact form matrix of \(h_i\) in the independent basis.
Compute a symmetric rational matrix \(\widetilde A\) by evaluating
each potential entry, once per unordered index pair, to error at most
\[
 \tau=\frac{\varepsilon_D\mu}{16r}.
\]
The kinetic matrix is exact. Since a symmetric matrix's operator
norm is at most \(r\) times its largest absolute entry,
\begin{equation}
 \|M^{-1/2}(\widetilde A-A)M^{-1/2}\|
 \le\varepsilon_D/16.
 \label{eq:wells-relative-matrix-error}
\end{equation}
The required entry precision has
\(O(D+K_M+\log r)\) bits and is therefore polynomial.
Let \(\widetilde\lambda\) be the minimum of the generalized
Rayleigh quotient for \((\widetilde A,M)\).
The variational characterization and
\eqref{eq:wells-relative-matrix-error} give
\[
 |\widetilde\lambda-\lambda_i^{\mathrm R}|
 \le\varepsilon_D/16.
\]

An initial rational bracket is \([-R_0,R_0]\), where
\[
 R_0=1+\frac{r\max_{a,b}|\widetilde A_{ab}|}{\mu}.
\]
Its bit length is polynomial. To bisect it at a rational \(z\),
test the inertia of \(\widetilde A-zM\) exactly.
For clarity, an elementary exact inertia algorithm repeatedly
uses congruence and Schur complements. A nonzero diagonal entry
is a one-dimensional pivot and contributes its sign. If every
remaining diagonal vanishes but an off-diagonal entry \(a\ne0\)
remains, pivot on the block
\(\left(\begin{smallmatrix}0&a\\a&0\end{smallmatrix}\right)\),
which contributes one positive and one negative direction.
A zero remaining matrix contributes only zero directions.
These operations give the full inertia. After denominators are
cleared, each Schur entry is a ratio of bordered minors; determinant
bounds again give polynomial bit growth.
Thus this test is a polynomial-time exact rational calculation,
including the case of zero eigenvalues.

Because \(M\) is positive definite,
\(\widetilde A-zM\) is positive semidefinite exactly when
\(z\le\widetilde\lambda\). Update the lower or upper endpoint
accordingly. Polynomially many bisections give a rational enclosure
\([\ell,u]\) of \(\widetilde\lambda\) of width at most
\(\varepsilon_D/8\). Arbitrary proximity to a bisection point incurs
no further precision requirement and no eigenvalue-gap assumption.

The fixed lower bound on \(D\) can be enlarged so that
\(Ce^{-1.16D}\le\varepsilon_D/8\); this is possible since
\((8/5)\log2<1.16\).
Combining the one-sided Ritz bias and the two-sided matrix error
gives the certified interval
\begin{equation}
 \lambda_i\in
 \left[\ell-\frac{\varepsilon_D}{16}
              -\frac{\varepsilon_D}{8},
       u+\frac{\varepsilon_D}{16}\right].
 \label{eq:wells-certified-enclosure}
\end{equation}
Its width is at most \(3\varepsilon_D/8\).
Notice that the Ritz bias is subtracted from the lower endpoint;
a Ritz upper bound alone would not suffice for tuning.
This proves the proposition.

All dimensions, operation counts, and working bit lengths are
polynomial in \(D,m,B\). The mesh and the independent subset may
change at different rational samples, but every such sample obeys
the same estimates. Constants controlling the polynomial degree,
the cutoff size, and the minimum admissible \(D\) are fixed once
from the uniform inequalities above. They are not input-dependent
advice or a search for an unknown spectral gap.
\end{proof}

\subsection{Simultaneous rational tuning}

The energy oracle now supplies certified intervals, while
\Cref{lem:wells-secants} supplies the monotonicity and weak coupling of
the tuning parameters. Combining these facts gives a rational algorithm:
each coordinate is bracketed separately, and simultaneous updates contract.
The stopping rule below uses an interval containing the target energy;
it never asks for the exact sign of an unknown eigenvalue difference.

\begin{proposition}[Tuning with fixed rounded charges]
\label{prop:tuning}
At any fixed rational centers satisfying the hypotheses of this
section, and with the integer charges in \eqref{eq:wells-charges}
already fixed, one can compute rational
\(\rho_1,\ldots,\rho_m\in[1,4]\) such that
\begin{equation}
 |\lambda_i+1/2|\le e^{-1.01D}\qquad(1\le i\le m).
 \label{eq:wells-tuned}
\end{equation}
The algorithm is deterministic and polynomial in \(D,m\) and the
input bit lengths. Every parameter it samples or outputs has
\(O(D)\) fractional bits.
\end{proposition}

\begin{proof}
For frozen parameters other than \(\rho_i\), the endpoint brackets,
continuity, and strict secant bound in \Cref{lem:wells-secants}
give a unique root of \(\lambda_i=-1/2\) in \([1,4]\).
Let \(T_i(\rho)\) be that root and let
\(T=(T_1,\ldots,T_m)\), using the old other coordinates in every
update. Combining \eqref{eq:wells-own-slope} and
\eqref{eq:wells-other-slope} gives
\begin{equation}
 \|T(\rho)-T(\widetilde\rho)\|_\infty
 \le\theta\|\rho-\widetilde\rho\|_\infty,\qquad
 \theta=Cm/D^2\le1/2.
 \label{eq:wells-contraction}
\end{equation}
Thus \(T\) has a unique fixed point \(\rho^*\) in the cube.
This also follows directly by iterating the map: its successive
differences are geometrically summable, and the limit is fixed
by continuity.

We implement each coordinate update with a certified scalar
bisection. Start with \([1,4]\). At its midpoint invoke
\Cref{prop:ritz-oracle}. If the whole resulting eigenvalue interval
lies below \(-1/2\), retain the upper half of the parameter interval;
if it lies above \(-1/2\), retain the lower half. If it contains
\(-1/2\), stop and return that sample. In this last case the true
energy differs from \(-1/2\) by at most \(\varepsilon_D\), so
\eqref{eq:wells-own-slope} puts the sample within
\(C\varepsilon_D/q'\) of the exact root.
Otherwise, after \(J=\lceil4D\rceil\) bisections return the midpoint
of the remaining parameter interval. All sign decisions were
certified, so that interval still contains the root and has length
\(3\,2^{-J}\). Since \(q'\le1\) and
\(2^{-4D}\ll\varepsilon_D\), the same update-error bound holds:
\begin{equation}
 \|\widehat T(\rho)-T(\rho)\|_\infty
 \le C\varepsilon_D/q'.
 \label{eq:wells-update-error}
\end{equation}
This comparison-or-stop rule never requests an exact sign of an
unknown irrational eigenvalue at a putative zero.

Start with all coordinates equal to two and perform \(J\) approximate simultaneous
updates. If \(e_j\) is the sup-norm distance to \(\rho^*\), then
\[
 e_{j+1}\le\tfrac12 e_j+C\varepsilon_D/q',
 \qquad
 e_J\le3\,2^{-J}+2C\varepsilon_D/q'.
\]
The own-coordinate upper slope and the other-coordinate variation
bound together give
\[
 |\lambda_i(\rho)-\lambda_i(\rho^*)|
 \le Cq'\|\rho-\rho^*\|_\infty.
\]
Consequently the output has
\[
 |\lambda_i+1/2|
 \le Cq'2^{-J}+C\varepsilon_D
 \le e^{-1.01D}
\]
for large enough \(D\), since
\((8/5)\log2>1.1>1.01\).

Each scalar bisection restarts from the rational interval \([1,4]\)
and makes at most \(O(D)\) midpoint operations. Its samples,
including early-return samples, therefore have \(O(D)\) fractional
bits independently of the number of simultaneous iterations.
There are \(O(mD^2)\) oracle calls. Their data bit lengths consist
of the fixed-center and rounded-charge bit lengths together with
these \(O(D)\) parameter bits, so \Cref{prop:ritz-oracle} proves
the total polynomial-time bound.
In particular the charge-rounding error is part of every evaluated
operator and is compensated by tuning; it is not introduced as an
uncontrolled perturbation afterward.
\end{proof}
